An object of this class represents a function in the application. A
\code{BPatch\_image} object (see description below) can
be used to retrieve a \code{BPatch\_function} object representing a
given function.

\begin{apient}
  std::string getName();
  std::string getDemangledName();
  std::string getMangledName();
  std::string getTypedName();
  void getNames(std::vector<std::string> &names);
  void getDemangledNames(std::vector<std::string> &names);
  void getMangledNames(std::vector<std::string> &names)
  void getTypedNames(std::vector<std::string> &names);
\end{apient}

\apidesc{
  Return name(s) of the function. The getName functions return the
  primary name; this is typically the first symbol we
  encounter while parsing the program; getName is an alias for
  getDemangledName. The getNames functions return all known
  names for the function, including any names specified by weak symbols.
}

\begin{apient}
  bool getAddressRange(Dyninst::Address &start, Dyninst::Address &end)
\end{apient}

\apidesc{
  Returns the bounds of the function; for non-contiguous functions,
  this is the lowest and highest address of code that
  the function includes.
}

\begin{apient}
  std::vector<BPatch_localVar *> *getParams()
\end{apient}

\apidesc{
  Return a vector of \code{BPatch\_localVar} snippets that refer to
  the parameters of this function. The position in the vector
  corresponds to the position in the parameter list (starting from
  zero). The returned local variables can be used to
  check the types of functions, and can be used in snippet expressions.
}

\begin{apient}
  BPatch_type *getReturnType()
\end{apient}

\apidesc{
  Return the type of the return value for this function.
}

\begin{apient}
  BPatch_variableExpr *getFunctionRef()
\end{apient}

\apidesc{
  For platforms with complex function pointers (e.g., 64-bit PPC)
  this constructs and returns the appropriate descriptor.
}

\begin{apient}
  std::vector<BPatch_localVar *> *getVars()
\end{apient}

\apidesc{
  Returns a vector of \code{BPatch\_localVar} objects that contain
  the local variables in this function. These
  can be used as parts of snippets in instrumentation. This function
  requires debug information to be
  present in the mutatee. If Dyninst was unable to find any local
  variables, this function will return an empty vector.
  It is up to the user to free the vector returned by this function.
}

\begin{apient}
  bool isInstrumentable()
\end{apient}

\apidesc{
  Return true if the function can be instrumented, and false if it
  cannot. Various conditions can cause a function to be
  uninstrumentable. For example, there exists a platform-specific
  minimum function size beyond which a function cannot
  be instrumented.
}

\begin{apient}
  bool isSharedLib()
\end{apient}

\apidesc{
  This function returns true if the function is defined in a shared library.
}

\begin{apient}
  BPatch_module *getModule()
\end{apient}

\apidesc{
  Return the module that contains this function. Depending on whether
  the program was compiled for debugging or the
  symbol table stripped, this information may not be available. This
  function returns NULL if module information was
  not found.
}

\begin{apient}
  char *getModuleName(char *name, int maxLen)
\end{apient}

\apidesc{
  Copies the name of the module that contains this function into the
  buffer pointed to by name. Copies at most maxLen
  characters and returns a pointer to name.
}

\begin{apient}
  enum BPatch_procedureLocation {
    BPatch_entry,
    BPatch_exit,
    BPatch_subroutine,
    BPatch_locInstruction,
    BPatch_locBasicBlockEntry,
    BPatch_locLoopEntry,
    BPatch_locLoopExit,
    BPatch_locLoopStartIter,
    BPatch_locLoopStartExit,
    BPatch_allLocations
  }
\end{apient}

\apidesc{}

\begin{apient}
  const std::vector<BPatch_point *> *findPoint(const BPatch_procedureLocation loc)
\end{apient}

\apidesc{
  Return the \code{BPatch\_point} or list of \code{BPatch\_points}
  associated with the procedure. It is used to select which type of
  points associated with the procedure will be returned.
  \code{BPatch\_entry} and \code{BPatch\_exit} request respectively the entry
  and exit points of the subroutine. \code{BPatch\_subroutine}
  returns the list of points where the procedure calls other
  procedures. If the lookup fails to locate any points of the
  requested type, NULL is returned.
}

\begin{apient}
  enum BPatch_opCode {
    BPatch_opLoad,
    BPatch_opStore,
    BPatch_opPrefetch
  }
\end{apient}

\apidesc{}

\begin{apient}
  std::vector<BPatch_point *> *findPoint(const std::set<BPatch_opCode> &ops)
  std::vector<BPatch_point *> *findPoint(const BPatch_Set<BPatch_opCode>& ops)
\end{apient}

\apidesc{
  Return the vector of \code{BPatch\_point}s corresponding to the set
  of machine instruction types described by the argument.
  Any combination of opcode type may be
  requested by passing an appropriate argument set containing the
  desired types. The instrumentation points created by
  this function have additional memory access information attached to
  them. This allows such points to be used for memory
  access specific snippets (e.g. effective address). The memory
  access information attached is described under Memory
  Access classes in section.
}

\begin{apient}
  BPatch_localVar *findLocalVar(const char *name)
\end{apient}

\apidesc{
  Search the function\'s local variable collection for name. This
  returns a pointer to the local
  variable if a match is found. This function returns NULL if it
  fails to find any variables.
}

\begin{apient}
  std::vector<BPatch_variableExpr *> *findVariable(const char *name)
  bool findVariable(const char *name, std::vector<BPatch_variableExpr> &vars)
\end{apient}

\apidesc{
  Return a set of variables matching name at the scope of this
  function. If no variables match in the local scope, then
  the global scope will be searched for matches. This function
  returns NULL if it fails to find any variables.
}

\begin{apient}
  BPatch_localVar *findLocalParam(const char *name)
\end{apient}

\apidesc{
  Search the function\'s parameters for a given name. A
  \code{BPatch\_localVar} pointer is returned if a
  match is found, and NULL is returned otherwise.
}

\begin{apient}
  void *getBaseAddr()
\end{apient}

\apidesc{
  Return the starting address of the function in the mutatee\'s address space.
}

\begin{apient}
  BPatch_flowGraph *getCFG()
\end{apient}

\apidesc{
  Return the control flow graph for the function, or NULL if this
  information is not available. The \code{BPatch\_flowGraph} is
  described in section.
}

\begin{apient}
  bool findOverlapping(std::vector<BPatch_function *> &funcs)
\end{apient}

\apidesc{
  Determine which functions overlap with the current function. Return
  true if other
  functions overlap the current function; the overlapping functions
  are added to the funcs vector. Return false if no
  other functions overlap the current function.
}

\begin{apient}
  bool addMods(std::set<StackMod *> mods)
\end{apient}

\apidesc{
  Implemented on x86 and x86-64

  Apply stack modifications in mods to the current function; the
  StackMod class is described in section. Perform error checking,
  handle stack alignment
  requirements, and generate any modifications
  required for cleanup at function exit. addMods atomically adds all
  modifications in mods; if any mod is found to be
  unsafe, none of the modifications in mods will be applied.

  addMods can only be used in binary rewriting mode.

  Returns false if the stack modifications are unsafe or if Dyninst
  is unable to perform the analysis required to
  guarantee safety.
}
